\section{Generative AI Use}\label{app:ai}

Two AI tools were used heavily for code development. Where possible,
queries are reproduced from session logs with minor typographical corrections
for readability but are otherwise verbatim; those corrections were themselves
made with Claude Sonnet 4.6 assistance.

\textbf{Claude Sonnet~4.6} (Anthropic), accessed via the Claude Code
command-line interface and web interface (\texttt{claude.ai/code}), was the
primary tool. All Claude-assisted commits carry a
\texttt{Co-Authored-By: Claude Sonnet~4.6} trailer that Claude Code
automatically appends, providing verifiable attribution in the git history.
The following queries and responses are reproduced from local Claude Code
session logs (\texttt{.jsonl} files stored by the CLI); the
priority-scheduler session (\S\ref{app:ai-scheduler}) was conducted via the
web interface and is documented from commit trailers only, as no local log
was retained.

\textbf{OpenAI Codex} was also used during development for feature
implementation, design discussion, and code-quality passes. Session logs were
not retained for Codex interactions and Codex does not automatically tag
commits, so specific queries cannot be reproduced.

\subsection*{A.2.1\quad Iterative Deepening Engine (6 May 2026)}

\textbf{Queries:}
\begin{quote}
  I want to implement an iterative deepening backwards search engine. The
  feature may not need to be a full engine, but the idea is to use the clap CLI
  args so the user can select via \texttt{--engine naive} or
  \texttt{--engine id}.
  I will be adding more optimisations later, but this is the starting point. The
  current architecture should support iterative deepening. Here is
  the algorithm:
  [Algorithm~2, p.~67, Hou 2021, supplied as full \LaTeX\ pseudocode].
\end{quote}

\begin{quote}
  I started an ultraplan session in the web browser and it began editing code.
  I would like to continue from the CLI instead, using this plan:
  [Plan: Iterative Deepening Backward Search Engine --- \texttt{--engine} flag,
    \texttt{ProofOptions::engine} field, ID loop calling
    \texttt{backwards\_search}
  at increasing depth limits, \texttt{engine/id.rs} module].
\end{quote}

\begin{quote}
  Please also make sure the engine can be configured in \texttt{config.toml},
  that everything is documented appropriately, and that all relevant functions
  have appropriate rustdocs.
\end{quote}

\textbf{Outcome:} Claude implemented the \texttt{--engine naive|id} CLI flag,
\texttt{ProofOptions::engine} field, and the iterative-deepening loop in a new
\texttt{engine/id.rs} submodule, adding rustdocs throughout
(commits \texttt{d16bb3c} and associated engine-modularisation commits).

\subsection*{A.2.2\quad SQLite Persistence (6 May 2026)}

\textbf{Queries:}
\begin{quote}
  I need to make the results from each run for each engine persist so I can
  analyse the data. How should I do this?
\end{quote}

\begin{quote}
  SQLite would be preferable --- can we use that instead?
\end{quote}

\begin{quote}
  The interface should be something like \texttt{--persist true} or
  \texttt{--persist false}, defaulting to \texttt{true}, with the database
  location configurable --- perhaps \texttt{--persist <DB\_FILE\_LOCATION>} with
  a default path specified in \texttt{config.toml}.
\end{quote}

\begin{quote}
  The CLI flag wiring should live in \texttt{cli/}, but the actual insertion
  logic should be in its own module such as \texttt{src/persistence}.
  What is your recommendation?
\end{quote}

\textbf{Outcome:} Claude designed and implemented the \texttt{src/persistence/}
module with an SQLite schema, \texttt{--persist} CLI flag with configurable
path, and a TDD test suite (commits \texttt{a8333b8}, \texttt{c5bc21f},
\texttt{5719e21}, \texttt{ba06ddf}, \texttt{2ca6f0e}, \texttt{a1de674}).

\subsection*{A.2.3\quad Six-Class Priority Scheduler (6 May
2026)}\label{app:ai-scheduler}

\textbf{Note:} This session was conducted via the Claude Code web interface;
no local session log is available. The session is documented from the
\texttt{Co-Authored-By: Claude Sonnet 4.6 <noreply@anthropic.com>} trailers
that Claude Code automatically appends to commits it helps author.

\textbf{Outcome:} Claude implemented the six-class LK$'$ rule priority scheduler
(closing rules; non-branching propositional; branching propositional;
  eigenvariable quantifiers; reusable quantifiers with visible witnesses;
reusable quantifiers with fresh fallback) and the \texttt{Priority} and
\texttt{PriorityId} engine variants
(commits \texttt{8d72635}, \texttt{663ac9a}, \texttt{d16bb3c}).

\subsection*{A.2.4\quad FOLIO Rust Integration (9 May 2026)}

\textbf{Queries:}
\begin{quote}
  I am happy to modify the prover so that it does not require a \texttt{.p}
  file --- it can be adapted to accept whatever format is needed.
  Please have a look at \texttt{theorem\_prover/} to determine the
  best approach.
\end{quote}

\begin{quote}
  Ensure every new function is appropriately documented with rustdocs, and
  maintain appropriate separation of concerns across modules, files, and
  functions. Create new submodules or files where needed to meet good software
  design standards.
\end{quote}

\textbf{Outcome:} Claude extended the grammar and parser to accept FOLIO
\texttt{.jsonl} problems directly without converting to TPTP \texttt{.p}
format, adding a new parser module with full rustdoc coverage.

\subsection*{A.2.5\quad Parallelisation and Channel-Based DB Writer
(9 May 2026)}

\textbf{Queries:}
\begin{quote}
  I would like to run problem files in parallel within
  \texttt{theorem\_prover/}.
  I am considering migrating persistence from SQLite to PostgreSQL to support
  this, though I am not sure that is the right approach.
\end{quote}

\begin{quote}
  For parallelisation: multiple files simultaneously within a single invocation.
\end{quote}

\begin{quote}
  Out-of-order progress output is acceptable. Please implement this with TDD.
  The timer logic can also be simplified, and the Ctrl-C handling can likely be
  simplified now that the owner thread can terminate the workers directly.
\end{quote}

\textbf{Outcome:} Claude implemented \texttt{rayon} \texttt{par\_iter} over
problem paths with an \texttt{AtomicUsize} progress counter, a channel-based
SQLite writer thread to avoid connection-sharing across threads, and an
exit-code fix (\texttt{process::exit} replaced with a return value to ensure
the writer thread flushes before exit)
(commits \texttt{c52a7bf}, \texttt{95cc45a}, \texttt{d48cc7b}).
